% ---------- Title block ----------
\thispagestyle{empty}
\begin{center}
  {\bfseries\LARGE 從頻道表現到留言社群洞察：\\[2pt]
   台灣 YouTube 頻道的 Audience Intelligence 分析\par}
  \vspace{6pt}
  {\bfseries\large 以 The DoDo Men－嘟嘟人為例\par}
  \vspace{10pt}
  {\small 社群媒體分析（Social Media Analytics）期末報告\par}
  \vspace{14pt}

  \renewcommand{\arraystretch}{1.25}
  \begin{tabular}{ll}
    \toprule
    \textbf{學號} & \textbf{姓名}\\
    \midrule
    B10704049 & 郭宜蓁\\
    R14921085 & 林業語\\
    B12705065 & 趙翊宏\\
    B11705037 & 鄭兆恩\\
    B11705042 & 趙仲文\\
    \bottomrule
  \end{tabular}
\end{center}

\vspace{8pt}
\begin{quotation}
\noindent\small\textbf{摘要.}\ 本研究將 YouTube 留言視為一種社群互動資料，而非影片的附屬回饋，
提出一套以「留言社群洞察」為核心的 \eng{audience intelligence} 分析框架。框架整合留言者活躍分層、
回訪率、\eng{co-commenter network} 社群偵測、情緒與 \eng{Aspect-Based Sentiment Analysis}（ABSA）
語意標註（採 Qwen 與 OpenAI API 雙模型投票加人工審核）以及 \eng{reply-thread conflict} 分析，並以 48 個
台灣頻道作為 \eng{broad benchmark}。
以台灣頻道 The DoDo Men－嘟嘟人為例，本研究分析 351 支影片、247{,}068 則主留言與 112{,}108 位留言者，
發現該頻道屬於「穩定討論引擎型」：回訪與核心黏著高於同儕、整體負面率偏低（7.5\%），但留言社群高度集中於三個
觀眾群，且少數影片出現被按讚放大的負面與回覆衝突。研究將數據結果轉化為內容策略、分眾溝通策略與炎上監測方向。
\end{quotation}

\vspace{4pt}
\hrule
\vspace{6pt}
